
\newpage
\section{Revisão bibliográfica}
\label{sc:revisao}

Esta seção sintetiza o referencial teórico e tecnológico que fundamenta o
Dança~AI em cinco eixos: dança de salão e forró universitário, visão
computacional e estimativa de pose, computação móvel e inferência
\textit{on-device}, inteligência artificial para análise de movimentos, e
trabalhos correlatos.

\subsection{Dança de Salão e Forró Universitário}
\label{sub:danca}

A dança de salão produz benefícios documentados à saúde física, à percepção
corporal e à integração social: Fonseca, Vecchi e Gama~\cite{fonseca2012}
distinguem as dimensões de esquema corporal e imagem corporal aprimoradas pela
prática; Rodrigues et al.~\cite{rodrigues2011} registram melhora significativa
de satisfação corporal em iniciantes; Paiva, Loures e Marinho~\cite{paiva2019danca}
relatam risco relativo de quedas de 0{,}2 em idosos praticantes; e Havt et
al.~\cite{havt2026dancar} identificam coordenação motora, socialização e
autoaceitação como motivações centrais. O forró universitário emergiu na década
de 1990 ao incorporar vocabulário da dança de salão --- giros, conexão corporal
e transferência de peso --- ao forró pé de serra~\cite{silva2022aspectos}, criando
um estilo com complexidade biomecânica adequada à análise computacional. A
sistematização pedagógica dessas modalidades ainda é escassa na
literatura~\cite{porto2023ressignificando}, reforçando a pertinência de uma
abordagem instrumentada.

A avaliação biomecânica dos movimentos baseia-se no ângulo articular $\theta$
calculado a partir de três \textit{landmarks} consecutivos --- proximal ($P$),
articular ($A$) e distal ($D$) --- segundo:

\begin{equation}
    \theta = \arccos\!\left(\frac{\overrightarrow{AP} \cdot \overrightarrow{AD}}
    {\left|\overrightarrow{AP}\right|\left|\overrightarrow{AD}\right|}\right)
    \label{eq:angulo}
\end{equation}

\noindent onde $\overrightarrow{AP} = P - A$ e $\overrightarrow{AD} = D - A$.
Aplicado quadro a quadro sobre os \textit{landmarks} fornecidos pelo MediaPipe,
esse cálculo produz séries temporais de ângulos que refletem a execução
dinâmica do movimento. Do ponto de vista rítmico, o forró universitário opera
em compassos binários com BPM entre 100 e 180; a sincronia sensoriomotora
com metrônomos é um mecanismo estabelecido para aprimorar a consistência
temporal do movimento humano~\cite{repp2013sms, karageorghis2019dance,
bood2013running, whitton2023sms}.

\subsection{Visão Computacional e Estimativa de Pose}
\label{sub:visao}

A estimativa de pose por redes neurais profundas tornou-se o método padrão
para análise de movimento em tempo real; Latyshev et
al.~\cite{latyshev2024computer} confirmam que arquiteturas otimizadas para
inferência em borda oferecem o melhor equilíbrio entre acurácia e viabilidade
operacional. O MediaPipe~\cite{lugaresi2019mediapipe} implementa a arquitetura
BlazePose~\cite{bazarevsky2020blazeposeondevicerealtimebody}, que opera em dois
estágios sequenciais --- detector de pessoa e estimador de pose sobre a região
de interesse --- produzindo 33~\textit{landmarks} tridimensionais com
\textit{visibility scores} a mais de 30~FPS em dispositivos \textit{mid-range}.
A validade das medições angulares a partir desses \textit{landmarks} foi
confirmada por Latreche et al.~\cite{LATRECHE2023112826} em comparação com
goniômetros clínicos. Os principais desafios em contextos de dança ---
auto-oclusão em giros, variações de ponto de vista e iluminação variável ---
são documentados por Tharatipyakul et al.~\cite{THARATIPYAKUL2024e36589} em
revisão de 45 trabalhos; o \textit{visibility score} do BlazePose é o
mecanismo adotado para suspender cálculos sobre \textit{landmarks} não
confiáveis~\cite{bora2023real}.

\subsection{Computação Móvel e Inferência \textit{On-Device}}
\label{subsub:otimizacao}

A inferência totalmente \textit{on-device} é determinada por três fatores:
(i)~\textbf{latência} --- o \textit{feedback} deve ser entregue em menos de
100~ms, janela incompatível com requisições de rede em condições reais de
conectividade móvel~\cite{li2020edgeai}; (ii)~\textbf{privacidade} ---
gravações contínuas do usuário não devem ser transmitidas a servidores
externos; (iii)~\textbf{disponibilidade} --- o aplicativo deve funcionar sem
conectividade. O TensorFlow Lite (TFLite)~\cite{lee2019ondevice} executa a
inferência com \textit{delegate} GPU (OpenGL ES~3.1) e NNAPI, com aceleração
média de 2--9$\times$ sobre CPU; a quantização INT8 reduz o tamanho do modelo
em 75\% sem degradação crítica de acurácia, viabilizando modelos embarcados em
\textit{mid-range}. O orçamento de latência distribui-se pelos estágios do
\textit{pipeline}: captura ($\approx$10~ms), inferência BlazePose
($\approx$20--40~ms), classificadores e módulo de ritmo ($\approx$10~ms), e
renderização ($\approx$16~ms), somando confortavelmente menos de 100~ms em
\textit{hardware} representativo~\cite{zhang2020mobipose}.

\subsection{Inteligência Artificial para Análise de Movimentos}
\label{sub:ia}

A classificação de movimentos de dança exige arquiteturas capazes de modelar
dependências temporais: arquiteturas CNN-LSTM combinam extração de
características locais por camadas convolucionais com modelagem de dependências
globais em sequências de \textit{landmarks}~\cite{ruiz2023dl_timeseries,
ncbi_dance_emotional}. Para os módulos de postura e identificação de movimentos
do Dança~AI, classificadores TFLite são treinados sobre coordenadas $xyz$ dos
\textit{landmarks} extraídas pelo MediaPipe --- o classificador de postura opera
sobre valores instantâneos de um único quadro; o classificador de movimentos opera
sobre séries temporais de quadros consecutivos. O dataset final consiste em
coordenadas, não em vídeos brutos, o que além de
preservar a privacidade dos participantes reduz a dimensionalidade da entrada
e elimina variações irrelevantes de iluminação, vestuário e fundo, tornando a
análise mais simples e precisa do que sobre a imagem completa. O classificador de postura é multi-label, avaliando cinco desvios de forma
independente --- ombro esquerdo caído, ombro direito caído, ombros encurvados,
inclinação do tronco para esquerda e inclinação do tronco para direita ---
podendo detectar múltiplos desvios simultaneamente; o classificador de
movimentos tem escopo restrito à identificação do passo em execução entre os
seis do forró, com a avaliação de qualidade da execução delegada ao módulo
DTW~\cite{bursa2023tflite, roggio2024pose}.

A comparação de sequências emprega o \textit{Dynamic Time Warping} (DTW), que
alinha duas séries temporais de comprimentos distintos minimizando a distância
acumulada entre elementos correspondentes, produzindo uma métrica invariante a
distorções temporais locais~\cite{dynamic_dance_warping}. A restrição de janela
de Sakoe-Chiba reduz a complexidade de $O(nm)$ para $O(n \cdot w)$ e
demonstrou acurácia superior ao DTW irrestrito em reconhecimento
gestual~\cite{10.1145/3598151.3598181, gesture_dtw_skeleton}. A eficácia do
\textit{feedback} visual depende da latência: atrasos superiores a 200~ms
comprometem a associação sensoriomotora~\cite{plos_feedback_delay,
pmc_movement_feedback_assist}; a sobreposição de esqueleto com codificação
cromática é a abordagem consolidada na
literatura~\cite{THARATIPYAKUL2024e36589}, com ganho empírico de acurácia de
55,9\% para 66,8\% em experimento com praticantes~\cite{stanford_ai_dance_coaching}.

\subsection{Trabalhos Correlatos}
\label{sub:correlatos}

A Tabela~\ref{tab:correlatos} compara o Dança~AI com os sistemas mais próximos
identificados na revisão. Nenhum dos trabalhos analisados reúne simultaneamente
inferência \textit{on-device} em \textit{smartphone} de consumo, avaliação
biomecânica por ângulos articulares, DTW para comparação de sequências e foco
em modalidade específica de dança de salão.

\begin{table}[h]
\centering
\caption{Comparação entre o Dança~AI e trabalhos correlatos}
\label{tab:correlatos}
\begin{tabular}{p{2.6cm}p{2.6cm}p{3cm}p{1.6cm}p{3.4cm}}
\hline
\textbf{Trabalho} & \textbf{Plataforma/ Execução} & \textbf{Foco de Avaliação} & \textbf{Feedback em Tempo Real} & \textbf{Limitação frente ao Dança~AI} \\
\hline
TransCNN-DSSS~\cite{Qu2024} & Servidor (GPU) & Qualidade do movimento (QSD: acurácia, fluidez, expressividade) & Não & Arquitetura pesada, sem viabilidade mobile \\
DanceFormer~\cite{danceformer} & Edge server dedicado & Análise de dança em tempo real & Sim (35,2~ms) & Depende de infraestrutura de borda externa \\
Coaching baseado em pose~\cite{stanford_ai_dance_coaching} & Protótipo de pesquisa & Acurácia de execução motora & Sim & Sem empacotamento como app autônomo \\
MobiPose~\cite{zhang2020mobipose} & Smartphone (on-device) & Estimativa de pose multipessoa & Sim & Sem avaliação biomecânica ou de qualidade \\
V-VIPE~\cite{levy2024v} & Não especificado & Embeddings de pose invariantes à vista & Não & Componente isolado, sem pipeline completo \\
Dynamic Dance Warping~\cite{dynamic_dance_warping} & Offline & Similaridade de sequências via DTW & Não & Algoritmo de comparação, sem integração mobile \\
\textbf{Dança~AI (proposto)} & \textbf{Smartphone (on-device)} & \textbf{Ângulos articulares + DTW} & \textbf{Sim (<100~ms)} & --- \\
\hline
\end{tabular}
\end{table}
